\section{An exactly central sparse LP and its oracle separations}\label{sec:lp}
We show that the spectral family arises as the normal matrix of the primal--dual Newton system of a small linear program. The state task asks for the normalized dual Newton direction as a quantum state; the algorithm may measure and discard auxiliary registers. The conversion task asks for a reusable coherent block encoding of $I-H_t$ with normalization one. Only the state task also receives an oracle that prepares the right-hand side.

\subsection{The linear program and its Newton equation}
Fix public parameters $\rho>1$ and $\delta>0$ with $\delta<1/8$ and $\rho\delta<1/4$. The parameter $t\in[\delta,\rho\delta]$ is hidden. Define
\begin{align*}
 u_\pm&=(1,\pm1)^T/\sqrt2,& q_+&=(1,1,2)^T/\sqrt6,\\
 v_0&=(1,-1,0)^T/\sqrt2,& k&=(1,1,-1)^T/\sqrt3,\\
 \sin\theta&=2\sqrt{2\delta},\quad 0<\theta<\pi/2,&
 q_-&=\cos\theta\,v_0+\sin\theta\,k,\\
 &&q_0&=-\sin\theta\,v_0+\cos\theta\,k.
\end{align*}
The vectors $q_+,q_-,q_0$ are orthonormal, and so are $u_+,u_-$. Put
\begin{equation}\label{eq:lp-family}
 A_t=u_+q_+^T+\sqrt t\,u_-q_-^T,\qquad
 H_t=A_tA_t^T=P_++tP_-,\qquad P_\pm=u_\pm u_\pm^T.
\end{equation}
Thus $\norm{A_t}=\norm{H_t}=1$, $A_t$ has rank two, and $\kappa(H_t)=1/t=\Theta_\rho(\delta^{-1})$. Every row of $A_t$ has at most three nonzero entries, every column at most two, and every entry has absolute value at most one. Sparsity here refers to these constant counts in a fixed-size example, not to growing dimension.

\begin{proposition}[An exactly central Newton system]\label{prop:lp-center}
The three-variable, two-equality linear program
\begin{equation}\label{eq:lp-primal}
 \min\ \mathbf1^Tx\qquad\text{subject to}\qquad
 A_tx=b_t:=A_t\mathbf1,\quad x\geq0
\end{equation}
has a nontrivial compact feasible segment on which the objective is not constant. The point $x=s=\mathbf1$, $y=0$ is exactly central with barrier parameter $\mu=1$. The normal equation of its affine-scaling predictor step and its solution are
\begin{align}
 H_t\Delta y&=b_t=a\bigl(u_++\sqrt{\delta t}\,u_-\bigr),
 &a&=2\sqrt{2/3},\label{eq:lp-normal}\\
 \Delta y&=a\bigl(u_++\sqrt{\delta/t}\,u_-\bigr).
 \label{eq:lp-direction}
\end{align}
\end{proposition}
\begin{proof}
The feasible affine line is $\mathbf1+\lambda q_0$. Since $q_0$ is nonzero and orthogonal to the strictly positive vector $q_+$, it has both positive and negative entries. The inequalities $1+\lambda(q_0)_j\geq0$ therefore bound $\lambda$ on both sides and admit an open neighborhood of zero. Their intersection is a nontrivial compact interval. Moreover,
\[
 \mathbf1^Tq_0=\cos\theta/\sqrt3>0,
\]
so the objective varies strictly along this line.
For the dual convention $A_t^Ty+s=\mathbf1$, centrality is the system
$A_tx=b_t$, $A_t^Ty+s=\mathbf1$, $x_js_j=\mu$, $x,s>0$.
The stated point satisfies it at $\mu=1$. Linearizing with complementarity target zero gives exactly
\[
 A_t\Delta x=0,\qquad A_t^T\Delta y+\Delta s=0,
 \qquad\Delta x+\Delta s=-\mathbf1.
\]
Eliminating $\Delta x,\Delta s$ yields $H_t\Delta y=A_t\mathbf1$. Finally,
$q_+^T\mathbf1=a$ and $q_-^T\mathbf1=a\sqrt\delta$, proving both formulas.
\end{proof}
The parameter $t$ changes how the equality constraints are written but not the feasible segment, since the row space of $A_t$ is fixed. Accordingly, the primal direction $\Delta x=-(\cos\theta/\sqrt3)q_0$ is independent of $t$, whereas the dual direction \eqref{eq:lp-direction} is not. The state lower bound therefore concerns the dual direction; the primal direction is public.

\subsection{Oracle specifications}
The following are public: the objective, the central point, the formulas above, and all spaces and basis changes below, including their dependence on $\delta$. No numerical value of $t$, $A_t$, $b_t$, or any coefficient that depends on $t$ is given as classical data; in particular, $\norm{b_t}=a\sqrt{1+\delta t}$ is not given, since it would reveal $t$. Each call to a supplied oracle or its inverse, controlled or not, costs one query, also when the algorithm chooses coherently which of these calls to make. As in the rest of the paper, oracle-independent gates and classical computation are free.

In the normal-matrix model, the oracle is the following unitary on $\C^2\otimes\C^2$, whose first factor is the block flag:
\begin{equation}\label{eq:lp-H-oracle}
 U_H(t)=\begin{pmatrix}
 H_t&\sqrt{I-H_t^2}\\
 \sqrt{I-H_t^2}&-H_t
 \end{pmatrix}.
\end{equation}
Its compression to flag $|0\rangle$ is exactly $H_t$. On $\C^2\otimes\operatorname{span}\{u_+\}$ and $\C^2\otimes\operatorname{span}\{u_-\}$ it acts as $R(1)$ and $R(t)$, respectively, where
\[
 R(z)=\begin{pmatrix}z&\sqrt{1-z^2}\\\sqrt{1-z^2}&-z\end{pmatrix}.
\]
In the factor model, where the oracle encodes the rectangular factor $A_t$, let $\mathcal R=\C^2$ be the row-index (output) space and $\mathcal C=\C^3$ the column-index (input) space; the oracle is the unitary on $\mathcal R\oplus\mathcal C$
\begin{equation}\label{eq:lp-A-oracle}
 U_A(t)=\begin{pmatrix}
 \sqrt{I_{\mathcal R}-A_tA_t^T}&A_t\\
 A_t^T&-\sqrt{I_{\mathcal C}-A_t^TA_t}
 \end{pmatrix}.
\end{equation}
Its off-diagonal block from $\mathcal C$ to $\mathcal R$ is $A_t$: for the canonical embeddings $J_{\mathcal R},J_{\mathcal C}$,
$J_{\mathcal R}^*U_A(t)J_{\mathcal C}=A_t$.
On the span of each ordered pair $(u_j,q_j)$ it acts as
$\left(\begin{smallmatrix}\sqrt{1-\sigma_j^2}&\sigma_j\\\sigma_j&-\sqrt{1-\sigma_j^2}\end{smallmatrix}\right)$,
where $\sigma_+=1$ and $\sigma_-=\sqrt t$; it acts as $-1$ on $q_0$. These blocks show that $U_A(t)$ is unitary and determine it on the whole space. Any identity padding needed for a qubit implementation is public and independent of $t$.

For the state task, both models also supply the unitary $B_t$ on $\mathcal R$ that prepares the normalized right-hand side; in the ordered basis $(u_+,u_-)$,
\begin{equation}\label{eq:lp-B-oracle}
 B_t=\frac1{\sqrt{1+\delta t}}
 \begin{pmatrix}1&-\sqrt{\delta t}\\\sqrt{\delta t}&1\end{pmatrix},
 \qquad B_tu_+=b_t/\norm{b_t}.
\end{equation}
We specify every block because unspecified blocks could reveal $t$ through unused registers.

\subsection{Matching query bounds for the Newton state}
Write $D(\omega,\tau)=\tfrac12\norm{\omega-\tau}_1$ for the trace distance and define
\begin{equation}\label{eq:lp-target}
 |\psi_t\rangle=
 \frac{u_++\sqrt{\delta/t}\,u_-}{\sqrt{1+\delta/t}},\qquad
 d_\rho=\frac{\sqrt\rho-1}{\sqrt{2(\rho+1)}}.
\end{equation}
By \eqref{eq:lp-direction}, $|\psi_t\rangle$ is the normalized dual Newton direction. The state task is to output a density operator $\omega_t$ with
$D(\omega_t,|\psi_t\rangle\langle\psi_t|)\leq\epsilon_0$ for every $t\in[\delta,\rho\delta]$, using a bounded worst-case number of queries. Here $\omega_t$ averages over all measurement outcomes; it is not conditioned on a favorable one.

\begin{theorem}[Preparing the Newton state with normal-matrix or factor access]\label{thm:lp-state}
Fix $\rho>1$ and $0<\epsilon_0<d_\rho/4$. In the normal-matrix and factor models, the minimum worst-case numbers of queries for the state task, counting calls to $B_t$, $B_t^*$, and the matrix or factor oracle, are
\begin{equation}\label{eq:lp-state-complexity}
 Q_H^{\rm state}(\delta,\epsilon_0)
 =\Theta_{\rho,\epsilon_0}(\delta^{-1}),\qquad
 Q_A^{\rm state}(\delta,\epsilon_0)
 =\Theta_{\rho,\epsilon_0}(\delta^{-1/2}).
\end{equation}
Since $\kappa(H_t)=\Theta_\rho(\delta^{-1})$ uniformly on the family, these orders are $\kappa(H_t)$ and $\sqrt{\kappa(H_t)}$ up to $\rho$-dependent constants.
\end{theorem}
\begin{proof}
For the lower bounds, the two endpoint target states have distance exactly $d_\rho$: their overlap is $(\sqrt\rho+1)/\sqrt{2(\rho+1)}$. Correct outputs for the two endpoints therefore have distance at least $d_\rho-2\epsilon_0>d_\rho/2$.
For $0\leq x,y\leq1/2$,
\[
 \norm{R(x)-R(y)}
 =2\left|\sin\frac{\arcsin x-\arcsin y}{2}\right|
 \leq2|x-y|.
\]
After a fixed change of basis, the $2\times2$ blocks of $U_A(t)$ satisfy the same bound in the singular value $\sigma_j$. At the endpoints, therefore,
\begin{align*}
 \norm{U_H(\rho\delta)-U_H(\delta)}&\leq2(\rho-1)\delta,\\
 \norm{U_A(\rho\delta)-U_A(\delta)}&\leq2(\sqrt\rho-1)\sqrt\delta.
\end{align*}
The preparation unitary is a rotation through angle $\arctan\sqrt{\delta t}$, so
\[
 \norm{B_{\rho\delta}-B_\delta}
 \leq(\sqrt\rho-1)\delta.
\]
These bounds are unchanged for inverses and controlled calls. Purify the algorithm, defer its measurements, and replace the queries of one endpoint computation by those of the other, one at a time. The distance between the final state vectors is at most the sum of the operator-norm distances of the replaced queries. Discarding registers does not increase trace distance, which is at most this vector distance. Consequently $T$ total queries give output distance at most $C_\rho T\delta$ in the normal-matrix model and $C_\rho T\sqrt\delta$ in the factor model. Comparing with $d_\rho/2$ proves both lower bounds, also for algorithms that call $B_t$.

For the upper bounds we estimate the single hidden parameter. Standard amplitude estimation \cite[Theorem 12]{brassard2002}, applied to a unitary preparation with success probability $a_0$, uses $O(M)$ calls to that preparation and its inverse and returns $\widehat a\in[0,1]$ such that
\begin{equation}\label{eq:lp-BHMT}
 |\widehat a-a_0|\leq
 \frac{2\pi\sqrt{a_0(1-a_0)}}{M}+\frac{\pi^2}{M^2}
\end{equation}
with probability at least $8/\pi^2>1/2$. Reflections about the public initial states and success subspaces are free; each reflection about the prepared state uses one counted call to the preparation and one to its inverse.
In the normal-matrix model, apply $U_H(t)$ to $|0\rangle u_-$ and count flag value $0$ as success; then $a_0=t^2$. Take $M\geq C/\delta$ and define $\widehat t$ by clipping $\sqrt{\widehat a}$ to $[\delta,\rho\delta]$. When \eqref{eq:lp-BHMT} holds, clipping does not increase the distance to $t$, and
\[
 |\widehat t-t|\leq
 \frac{|\widehat a-t^2|}{t}
 \leq\delta\left(\frac{2\pi}{C}+\frac{\pi^2}{C^2}\right).
\]
In the factor model, apply $U_A(t)$ to $J_{\mathcal C}q_-$ and declare the output subspace $\mathcal R$ successful; then $a_0=t$. Taking $M\geq C/\sqrt\delta$ and clipping $\widehat a$ instead gives
\[
 |\widehat t-t|\leq
 \delta\left(\frac{2\pi\sqrt\rho}{C}+\frac{\pi^2}{C^2}\right)
\]
on the same success event. Choose $C=C_{\rho,\epsilon_0}$ so that these bounds are at most $\epsilon_0\delta$. Taking the median of a fixed odd number of independent repetitions makes the failure probability at most $\epsilon_0/2$; the number of repetitions depends only on $\epsilon_0$.

Given the resulting estimate, prepare $|\psi_{\widehat t}\rangle$ using a rotation in the known basis. Its angle is $\varphi(t)=\arctan\sqrt{\delta/t}$, with
\[
 |\varphi'(t)|=\frac{\sqrt\delta}{2\sqrt t(t+\delta)}
 \leq\frac1{4\delta}\qquad(t\geq\delta).
\]
On successful estimation, this pure state is thus within trace distance $\epsilon_0/4$ of $|\psi_t\rangle$. Convexity of trace distance bounds the error of the average output by $\epsilon_0/4+\epsilon_0/2<\epsilon_0$, including all failure events. These algorithms make no calls to $B_t$ and use $O_{\rho,\epsilon_0}(\delta^{-1})$ or $O_{\rho,\epsilon_0}(\delta^{-1/2})$ matrix or factor queries, respectively. Rounding $M$ upward to a power of two changes only a constant factor. This proves the matching upper bounds.
\end{proof}
The identity $\Delta y=(A_t^T)^+\mathbf1$ also follows directly from the singular vectors, and the squared overlap of $\mathbf1/\sqrt3$ with $\operatorname{span}\{q_+,q_-\}$ is $(8/9)(1+\delta)$, bounded away from zero. Improvements from factor access, and the role of this overlap, are established in \citet[Sections 1.4 and 5.4--5.5]{orsucci2021}. Here the explicit estimation argument avoids the overhead of a generic solver and gives the exact orders at fixed output error. The algorithm outputs a state; it is not a reusable coherent converter.

\subsection{Converting to the complement on the LP family}
Let $Q_H^{\rm LP}(\delta,\eta)$ be the minimum number of queries of a single converter whose unitary has, for every $t\in[\delta,\rho\delta]$, a designated block within $\eta$ of
\[
 I-H_t=(1-t)P_-.
\]
For this task the only input oracle is \eqref{eq:lp-H-oracle}; $B_t$ is not supplied. All other data remain public. Define $Q_A^{\rm LP}$ analogously, with \eqref{eq:lp-A-oracle} as the only oracle and output on $\mathcal R$.

\begin{theorem}[Complement costs on the LP family]\label{thm:lp-compiler}
Fix $K>0$ and let $\ell=\min\{j\geq0:G_j(\rho)\leq K\}$. For sufficiently small $\delta$, a converter that uses only the normal-matrix oracle has query complexity
\begin{equation}\label{eq:lp-staircase}
 Q_H^{\rm LP}(\delta,K\delta)=
 \begin{cases}
 0,&\ell=0,\\
 \Theta_{\rho,K,\ell}(\delta^{-1+1/(2\ell)}),&\ell\geq1,
 \end{cases}
\end{equation}
where $G_j(\rho)$ are the thresholds defined in \eqref{eq:G-definition}; as in Theorem~\ref{thm:staircase}, $K=G_j(\rho)$ belongs to the cheaper tier. For every fixed $\beta>0$, uniformly for $0<\eta\leq\delta^{1+\beta}$,
\begin{equation}\label{eq:lp-joint}
 Q_H^{\rm LP}(\delta,\eta)
 =\Theta_{\rho,\beta}\bigl(\delta^{-1}\log(1/\eta)\bigr)
 =\Theta_{\rho,\beta}\bigl(\kappa(H_t)\log(1/\eta)\bigr).
\end{equation}
In contrast, with factor access, two queries implement $I-A_tA_t^T$ exactly with normalization one.
\end{theorem}
\begin{proof}
In the public eigenbasis, the normal-matrix oracle is the direct sum of the fixed block $R(1)$ and the block $R(t)$. Replace the latter by
$\left(\begin{smallmatrix}\sin\tau&\cos\tau\\\cos\tau&-\sin\tau\end{smallmatrix}\right)$ with a real parameter $\tau$; at $\tau=\arcsin t$ this is the actual oracle. Every matrix element of a $T$-query converter is now a bounded trigonometric polynomial of degree at most $T$, even if its gates mix the two eigenspaces. Its designated $u_-$ diagonal entry approximates $1-\sin\tau$ on the promised interval. Thus every lower bound in Sections~\ref{sec:fixed}--\ref{sec:joint} that uses correctness only on the low band applies without change. In particular, Theorem~\ref{thm:hierarchy} gives the lower bounds for all $\ell\geq1$, and Theorem~\ref{thm:joint-exterior} (or Theorem~\ref{thm:FR-all}) gives the high-accuracy lower bound. The reduction does not cover circuits that also call $B_t$, because the entries of $B_t$ would not obey the same degree induction.

The polynomial upper-bound constructions apply to $H_t$ with the single-point high band $c=1$. They prove the upper bounds for every $\ell\geq1$ and \eqref{eq:lp-joint}, using $K=\eta/\delta$ and
$\log(\delta/\eta)=\Theta_\beta(\log(1/\eta))$ in the stated regime. For the coarse tier, the public operator
$(1-m_\rho\delta)P_-$ is a contraction and differs from $(1-t)P_-$ by at most $G_0\delta$. Its standard unitary dilation uses no oracle queries. In the general $c=1$ problem of Theorem~\ref{thm:staircase}, the coarse tier costs a nonzero constant number of queries because the spectral projectors are not given; here the eigenspaces are public.

For factor access, the real even polynomial $p(s)=1-s^2$ has degree two and lies in $[0,1]$ on $[-1,1]$. The even, left singular-vector version of QSVT \cite[Corollary 18, extended version]{gilyen2019}, applied with the public row and column projectors in \eqref{eq:lp-A-oracle}, therefore gives exactly
\[
 \sum_{j=\pm}(1-\sigma_j^2)u_ju_j^T
 =I_{\mathcal R}-A_tA_t^T
\]
as a unit-normalized block using two queries. The output acts on the two-dimensional left singular-vector space $\mathcal R$; the extra column null vector does not create an extra output eigenvalue. Equivalently, if $\Pi_{\mathcal C}=J_{\mathcal C}J_{\mathcal C}^*$, then
\[
 J_{\mathcal R}^*U_A(t)(I-\Pi_{\mathcal C})U_A(t)^*J_{\mathcal R}
 =I_{\mathcal R}-A_tA_t^T.
\]
A free block encoding of the public projector $I-\Pi_{\mathcal C}$ makes this an explicit two-query construction.
\end{proof}
The same scalar-oracle reduction applies to the finite-index and growing-index bounds of Section~\ref{sec:joint}, together with their stated limitations at intermediate accuracy. In particular, error $o(\delta)$ requires $\Omega_{\rho,\varepsilon}(\delta^{-1+\varepsilon})$ queries for every fixed $\varepsilon>0$, and no finite number of normal-matrix oracle queries achieves zero error on the whole interval of $t$. If $t$ were public, Proposition~\ref{prop:two-point} would give an exact two-query conversion for the known spectrum $\{t,1\}$, even with unknown eigenspaces; with the public eigenspaces of this example, direct synthesis would need no queries. The lower bounds therefore depend on $t$ being hidden and on correctness for every $t\in[\delta,\rho\delta]$.

\subsection{Access assumptions and scope}
The separation disappears under digital access to matrix entries. For example,
\[
 (H_t)_{11}=(1+t)/2,\qquad
 (A_t)_{11}=\frac1{\sqrt2}
 \left(\frac1{\sqrt6}+\sqrt t\,(q_-)_1\right),\qquad
 (q_-)_1=\frac{\cos\theta}{\sqrt2}+\frac{\sin\theta}{\sqrt3}>0.
\]
A single exact query to either displayed entry reveals $t$, and a constant number of sufficiently precise digital queries approximates $t$ well enough for any fixed state error; precision and arithmetic costs are counted separately. Since the dimension and eigenbases are known, such an approximation allows direct state preparation or synthesis of a complement dilation. This is consistent with the digital sparse-access constructions of \citet[Section 4.3]{orsucci2021}; indeed $H_t$ is strictly diagonally dominant, with diagonal $(1+t)/2$ and off-diagonal $(1-t)/2$.

A linear combination of unitaries gives a one-query encoding with normalization two: preparing a control qubit in $|+\rangle$, applying $\operatorname{diag}(I,-U_H(t))$, and projecting the control onto $|+\rangle$ gives $(I-H_t)/2$. Only the requirement of normalization exactly one rules out this shortcut. Orsucci and Dunjko already identify a normalized encoding of $I-\eta A$ as a substantive input requirement \cite[Section 4.3]{orsucci2021}; our lower bounds quantify its cost from a plain matrix oracle when the hidden parameter ranges over an interval.

These results are query separations for the specified quantum (analog) oracles and output tasks. They are not lower bounds for all algorithms that solve the LP, nor end-to-end lower bounds for quantum interior-point methods. In particular, the primal predictor is public, a solver need not construct a unit-normalized complement, and the query count does not measure classical runtime or digital precision.
